\section{The design idea behind Dagents}
\label{sec:design}

Dagents is an open-source framework for analytics and machine learning across
data sources that keep their data~\citep{dagents2026repo}. It started as code
shared by other products, such as Watchdog, and the Python package
\code{dagents} was first published to PyPI as version 0.1.0a1, an alpha
release. This section states the idea behind its design as ten principles in
three groups. Each principle names the code that implements it, so that a
reader can check the claim.

\subsection{The parts}

\Cref{tab:layers} lists the layers. One rule decides where code goes:
decisions go in the OCaml planners, and anything with side effects goes in
Python or Java.

\begin{table}[h]
\centering
\small
\begin{tabularx}{\linewidth}{@{}l l Y Y@{}}
\toprule
Layer & Language & Does & Does not \\
\midrule
Planners & OCaml & Validate, plan, route and render: governance, federated
  rounds, pipeline order, model families, Kubernetes and Compose files &
  Open connections, train models, serve HTTP \\
Agents and services & Python & The local agent (LMA) at each source, the
  global agent (GMA), the core, pipeline and model services; data access,
  training, state & Contain decision rules \\
Orchestration & Java & Control and core APIs for JVM consumers & Machine
  learning \\
\bottomrule
\end{tabularx}
\caption{The three layers of Dagents and what each one owns.}
\label{tab:layers}
\end{table}

A \textbf{local monitoring agent} (LMA) runs at each data source, such as a
hospital, a database or a tenant. It profiles the local data, runs local jobs,
reports summaries, and enforces governance at three boundaries. A
\textbf{global monitoring agent} (GMA) registers the LMAs, plans federated
rounds, sends each site its job, combines results and decides, with a person,
whether a candidate model may be released. The two agents call each other only
through their HTTP APIs, with bearer tokens of fixed scopes.
\Cref{fig:dagents} shows where each decision is made and enforced.

\begin{figure}[t]
\centering
\begin{tikzpicture}[
  node distance=3mm and 4mm,
  box/.style={draw=dgInk, rounded corners=2pt, align=center, font=\footnotesize, inner sep=3pt, minimum height=7mm},
  store/.style={draw=dgInk, cylinder, shape border rotate=90, aspect=0.18, align=center, font=\scriptsize, minimum height=9mm, minimum width=12mm},
  gate/.style={draw=dgInk, fill=dgSoft, font=\scriptsize\ttfamily, inner sep=2pt},
  part/.style={draw=dgRule, rounded corners=3pt, inner sep=4pt},
  arr/.style={-{Stealth[length=1.8mm]}, draw=dgInk},
  ask/.style={-{Stealth[length=1.6mm]}, draw=dgRule, dashed},
  lbl/.style={font=\scriptsize\itshape, text=dgInk}
]
% one site
\node[store] (st) {site\\records};
\node[gate, right=4mm of st] (read) {before\_read};
\node[box, right=4mm of read] (job) {local job on\\guarded data};
\node[gate, above=3mm of job] (train) {before\_train};
\node[gate, right=4mm of job] (send) {before\_send};
\node[part, fit=(st)(read)(job)(train)(send), label={[lbl]below:Site \(i\): its LMA and Guard (one per source)}] (site) {};
% coordinator
\node[box, right=14mm of send] (gma) {GMA: round plan,\\readiness, registry};
\node[gate, below=3mm of gma] (rel) {before\_release};
\node[box, below=3mm of rel] (cand) {candidate \(\rightarrow\) gates\\\(\rightarrow\) person approves};
\node[part, fit=(gma)(rel)(cand), label={[lbl]below:Coordinator (GMA)}] (coord) {};
% planner
\node[box, below=20mm of job, minimum width=60mm, fill=dgSoft] (plan) {\code{dagentsc}: pure OCaml planners (governance, federation)};
\draw[arr] (st) -- (read);
\draw[arr] (read) -- (job);
\draw[arr] (train) -- (job);
\draw[arr] (job) -- (send);
\draw[arr] (send) -- node[lbl, above] {counts,} node[lbl, below] {metrics} (gma.west);
\draw[arr] (gma.north) to[bend right=30] node[lbl, above] {job + manifest digest} (train.north);
\draw[arr] (rel) -- (cand);
\draw[ask] (read.south) -- (read.south |- plan.north);
\draw[ask] (send.south) -- (send.south |- plan.north);
\draw[ask] (rel.west) -| ($(plan.north east)+(-3mm,0)$);
\node[lbl, below=1mm of plan] {dashed: the Guard asks for a plan (JSON on standard input); no plan means deny};
\end{tikzpicture}
\caption{Where Dagents decides and where it enforces. The LMA at each site
enforces three boundaries, the GMA enforces the fourth, and every Guard asks the
same planner. The clipped model update goes from \code{before\_send} to a
write-only update store in the GMA, which is not drawn.}
\label{fig:dagents}
\end{figure}

\subsection{Where computation runs}

\paragraph{P1. Compute at the source.} Each source runs its own LMA and keeps
its records in its own store. In the healthcare demo, each hospital has one
PostgreSQL schema and one role that cannot log in. Every database operation runs
in a single transaction that starts with \code{SET LOCAL ROLE} and checks
\code{current\_user}, so a query for one hospital cannot read another hospital's
rows. The coordinator's role has no grant on any hospital schema. What crosses
to the coordinator is listed in \cref{tab:crosses}.

\paragraph{P2. Delegate the learning protocol.} Dagents does not implement a
federated optimiser. The GMA runs rounds through a \code{FederationEngine}
interface; the in-process engine is a simulator for tests, and the HTTP engine
moves jobs to the LMAs and holds clipped updates until aggregation. How updates
are combined comes from a \emph{model kind} that the application contributes,
so no averaging rule lives in the framework. A specialist runtime such as
Flower or NVIDIA FLARE can sit behind the same interface. Dagents decides who
takes part and whether a result may be released.

\subsection{How decisions are made}

\paragraph{P3. Separate deciding from doing.} This is the GRAILS split. The
planner chooses a strategy for each requested field and never sees a record.
The Guard, in Python, applies the plan to the records, drops every field the
request did not ask for, and writes an audit record. A change of policy is a
change to a classification; the code that reads records does not change.

\paragraph{P4. Make every decision a total, typed function.} Sensitivity, trust,
granularity, strategy, decision, guard boundary and model status are algebraic
data types. The strategy table is one \code{match} over the triple
(sensitivity, trust, granularity), with 27 arms that cover all 45 combinations
and no wildcard (\cref{lst:select}). Adding a sensitivity level or a
granularity makes the match non-exhaustive. The OCaml compiler reports that as a
warning, and dune's default development profile, which the repository's CI
builds on every change, treats it as an error. So the change cannot pass CI
until every combination has a strategy. A build in dune's release profile would
only warn; we checked both behaviours on a small example.
The planner modules contain no input or output, no clock and no randomness, so
a plan can be computed again from its request during an audit. Where a decision
depends on time, such as whether a stored update is too old, the caller passes
the age in. Tests check three properties of the table: raising trust never
makes the protection stricter, a model update is never released unchanged, and
high-sensitivity data is refused to a low-trust requester at every granularity.
The table is also monotone in sensitivity; we checked this with a copy of the
table, and the repository has no test for it yet.

\begin{lstlisting}[language=OCaml, float=t, caption={Part of the strategy table in the governance planner. The full match has 27 arms and no wildcard.}, label={lst:select}]
let select_strategy sensitivity trust granularity minimum_cohort =
  match (sensitivity, trust, granularity) with
  | HighSensitivity, LowTrust, (CellGrain | RowGrain | ColumnGrain
                                | TableGrain | ModelUpdateGrain) -> Refuse
  | HighSensitivity, ModerateTrust, (CellGrain | RowGrain) -> Redact
  | HighSensitivity, ModerateTrust, (ColumnGrain | TableGrain) ->
      AggregateOnly minimum_cohort
  | HighSensitivity, ModerateTrust, ModelUpdateGrain -> AddNoise noise_scale
  (* ... 22 more arms ... *)
  (* An update is never released unprotected, because it is derived
     from records rather than being one. *)
  | LowSensitivity, HighTrust, ModelUpdateGrain -> ClipContribution clip_bound
\end{lstlisting}

\paragraph{P5. Put the planner behind a process boundary.} Services call the
planners by running the \code{dagentsc} program with JSON on standard input,
never through a foreign-function interface. Each call has limits: 10 seconds by
default, and 16 MiB for the input and for the output. The runner kills a planner
that exceeds them, together with its process group. An operator can pin the
SHA-256 of the planner binary; with another hash the runner starts nothing. A
separate process isolates failures, lets each side be upgraded on its own, and
means the decision logic can be replaced by any program that honours the same
JSON contract.

\paragraph{P6. Deny when there is no decision.} The Guard fails closed. If the
planner cannot be started, runs past its limits, prints something that is not a
plan, returns a strategy the Guard cannot apply, or if applying the plan to the
records fails, the Guard denies the request and records the denial. The
repository's rules forbid a fallback that permits a request when the planner
fails.

\paragraph{P7. Keep policy as declared data.} Field sensitivities, approved
purposes and the minimum cohort size are a \emph{data classification}, which an
application contributes through an extension or registers at runtime. A runtime
registration cannot replace a classification that an extension contributed (the
API answers \code{409}), and an agent started for one site refuses a request
whose inline classification differs from the registered one. A caller therefore
cannot choose the policy its own request is checked against. The framework
itself contains no domain words: nothing in the \code{dagents} package knows
what a stroke is. The healthcare demo's whole integration is one declarative
extension file of 136 lines.

\subsection{What decisions cover, and what they leave behind}

\paragraph{P8. Govern every boundary that data crosses, with one decision
model.} The same planner and Guard serve four boundaries:
\begin{itemize}
  \item \code{before\_train} (LMA): may the site train on the requested fields
        for this round? Checked when the job is offered.
  \item \code{before\_read} (LMA): which records and fields may the job read?
        The job runs on the guarded data, not on the raw records.
  \item \code{before\_send} (LMA): may the round's result, with its counts,
        metrics and bounded model update, leave the site for the number of
        records the round used?
  \item \code{before\_release} (GMA): may the candidate model be released?
\end{itemize}
To govern what a federated round sends, Dagents adds a fifth granularity to the
four of GRAILS: the \emph{model update}, whose strategies are clipping to a
norm bound and adding noise. A model update always gets some protection,
because it is derived from the records. The update vector never appears in a
site's result, because results are shown to the coordinator: it goes through
\code{before\_send}, which clips it, to a write-only update store. The result
carries only the update's norm and digest. Two request-level checks apply
before any per-field lookup: the declared purpose must be one the
classification approves, and a column, table or model-update request must cover
at least the minimum cohort. If either fails, every field is refused.

\paragraph{P9. Separate a candidate from a release.} Aggregation produces a
candidate model, never a release. A release gate whose metric is missing blocks
the release instead of passing. Every change of a model version's status
(candidate, released, rejected, superseded, with at most one released) comes
from the planner, and the model registry applies the answer without a status
rule of its own. A person's approval releases a candidate only when the stored
decision says release, the Guard permitted it, the metrics came from the
study's own rounds, and the model it was compared with is still the one in use.
The round controller also refuses a reused round identifier, a result from a
site the plan did not select, and a second aggregation of the same round.

\paragraph{P10. Leave evidence with every decision.} Every Guard decision,
permitted or not, writes an audit record. Each record holds the SHA-256 digest
of its own content and of the previous record, so editing or deleting a record
breaks every digest after it. Each record names the planner that decided: the
package version and the SHA-256 of the planner binary. A plan carries the
obligations the Guard must discharge and a filtering score in the sense of
GRAILS. A round manifest has a content digest, each excluded site has a recorded
reason, and noise is seeded from the request, so an auditor can reproduce every
value; \cref{sec:limitations} explains why this choice also weakens the noise. In the healthcare demo, every successful JSON answer lists the agent
calls behind it.

\subsection{What the principles add up to}

Taken together, the principles do two things. First, they move the decisions
that the rules in \cref{sec:introduction} call for into one place, the
planners, where they can be read as a table, tested exhaustively and computed
again from a stored request. Second, they move enforcement to where the data
is, and apply it at every boundary that data or something derived from it
crosses, not only at the response to a query. GRAILS showed the first for one
application; Dagents applies it to a study that spans several organisations.

\paragraph{What Dagents does not do.} It does not prove a differential privacy
bound: clipping and noise bound each update, but there is no privacy accountant
or tracked budget. Round digests detect changes but are not signatures. The
audit chain detects tampering in a test but is not an append-only audit store.
The trust score is computed from attributes the caller supplies; Dagents does
not verify identities. These limits are discussed in \cref{sec:limitations}.
